\section{The cumulative hierarchy and rank}
\label{sec:cumulative-hierarchy-02}

\begin{definition}[Cumulative hierarchy]
By transfinite recursion define
\[
  V_0=\varnothing,\qquad
  V_{\alpha+1}=\Pow(V_\alpha),\qquad
  V_\lambda=\bigcup_{\beta<\lambda}V_\beta
\]
when \(\lambda\) is a limit ordinal.  The symbol
\[
  V=\bigcup_{\alpha\in\operatorname{Ord}}V_\alpha
\]
is class notation, not a set-forming union; \(\operatorname{Ord}\)
denotes the class of all ordinals, which the Burali-Forti argument shows to
be a proper class.
\end{definition}

\begin{theorem}[Hierarchy theorem]
Assume \(\ZF\).  Then:
\begin{enumerate}
  \item every \(V_\alpha\) is transitive;
  \item if \(\alpha\leq\beta\), then \(V_\alpha\subseteq V_\beta\);
  \item every set belongs to some \(V_\alpha\); and
  \item for every set \(x\), there is a least ordinal \(\alpha\) such that
        \(x\subseteq V_\alpha\).
\end{enumerate}
\end{theorem}

\paragraph{Proof sketch.}
The first two clauses follow by transfinite induction.  For the third,
suppose some set \(x\) belongs to no stage.  Foundation, applied to the
set of members of the transitive closure of \(\set{x}\) that belong to no
stage, supplies a membership-minimal counterexample.
All of its members occur at stages; replacement bounds their stage indices by
a set of ordinals.  Their union gives one stage containing every member, so the
counterexample appears at the next stage.  The fourth clause uses ordinal
well-ordering.  The dependency structure is recorded; routine coding details
are omitted.

\begin{definition}[Rank]
The rank of a set \(x\) is
\[
  \rank(x)\defeq
  \min\set{\alpha:x\subseteq V_\alpha}.
\]
Equivalently,
\[
  \rank(x)=\sup\set{\rank(y)+1:y\in x},
\]
where, for a set \(S\) of ordinals, \(\sup S\defeq\bigcup S\) is the
least ordinal that is at least every member of \(S\).
The equivalence is a theorem obtained from the hierarchy theorem and
transfinite recursion.
\end{definition}

\begin{example}[The first stages]
\[
  V_0=\varnothing,\qquad
  V_1=\set{\varnothing},\qquad
  V_2=\set{\varnothing,\set{\varnothing}},
\]
and \(V_3=\Pow(V_2)\) has four members.  Consequently
\(\rank(\varnothing)=0\), \(\rank(\set{\varnothing})=1\), and
\(\rank(V_2)=2\).  The index of the first stage containing \(x\) as a member
is \(\rank(x)+1\); it is not \(\rank(x)\).
\end{example}

